% =====================================================================
%  ХОЁР. НӨХЦӨЛ БАЙДЛЫН ШИНЖИЛГЭЭ
% =====================================================================
\chapter{НӨХЦӨЛ БАЙДЛЫН ШИНЖИЛГЭЭ}

\section{Макро орчны шинжилгээ}

\subsection{Улс төрийн нөхцөл байдал, хууль эрхзүйн зохицуулалт}

Монгол Улсад мэдээллийн технологийн салбарыг дэмжих бодлого бүх түвшинд
баримтлагдаж байна. «Цахим орчны аюулгүй байдлын тухай», «Хувийн нууцын
тухай» болон «Арилжааны нууцын тухай» хуулиуд нь хэрэглэгчийн хувийн
мэдээллийг цуглуулах, хадгалах, боловсруулахад тавигдах шаардлагыг
зохицуулдаг. Мөн «Төлбөр, төлбөрийн системийн тухай» хууль болон Монголбанкны
дүрэм журмын дагуу цахим төлбөрийн үйлчилгээг зөвхөн тусгай зөвшөөрөл
бүхий этгээд (банк, төлбөрийн үйлчилгээний лицензтэй компани) гүйцэтгэх
үүрэгтэй тул систем нь өөрийн гэсэн төлбөрийн хэрэгсэл гаргахгүй,
батлагдсан төлбөрийн үйлчилгээ үзүүлэгчид (QPay, SocialPay, карт
процессинг) холбож ажиллана \cite{privacylaw}.

Дээрх зохицуулалтууд нь төслийн хэрэгжилтэд саад болохгүй, харин ч
цахим төлбөр, хувийн мэдээллийн хамгаалалтын чиглэлээр ажилладаг
байгууллагуудад тавигдах шаардлагыг хатуу мөрдүүлж, технологийн
шийдлийн эрэлт нэмэгдүүлэх хүчин зүйл болно.

\subsection{Хүн ам зүй, нийгэм соёлын хүчин зүйл}

Улаанбаатар хотын хүн амын 60 гаруй хувь нь 35-аас доош насны
идэвхтэй хэрэглэгчид бөгөөд гар утасны нэвтрэлт, цахим төлбөрийн
хэрэглээ нь он жилүүдэд тасралтгүй өсөж байна. Хөгжил дэгжилтийн
судалгаагаар гар утасны интернэт хэрэглэгчдийн дийлэнх нь нэг буюу
хэд хэдэн мобайл төлбөрийн аппликешнийг өдөр тутамдаа ашигладаг болохоор
цахим захиалга, төлбөрийн системд хэрэглэгчид дасан зохицох,
сэтгэл ханамжийн хүлээлт нэлээд өндөр байна.

Мөн «кафе-т цаг зав гаргах» нь залуучуудын амьдралын хэв маягийн нэг
хэсэг болж үүссэн нь кофе шопын салбарын өсөлт, үйлчилгээний чанарын
өрсөлдөөнийг дэмжинэ.

\subsection{Эдийн засгийн хүчин зүйл}

Сүүлийн жилүүдэд Монгол Улсын эдийн засгийн өсөлт тогтвортой байж,
хэрэглэгчдийн зардлын бүтцэд үйлчилгээний зардал эзлэх хувь нэмэгдэж
байна. Мөнгөн гүйлгээний биш, цахим гүйлгээний тооллого нэмэгдсээр
байгаа нь төлбөрийн интеграцитай системүүдийн ажиллагааг хөнгөвчилнө.
Гэхдээ төгрөгийн ханшийн хэлбэлзэл, импортын барааны үнийн өсөлт нь
кофе шопын өртгийн бүтцэд дарамт үүсгэж байгаа тул системийн үнийн
бодлого нь салбарын өртгийн бүтцэд тохирсон, уян хатан байх шаардлага
гарч байна.

\section{Микро орчны шинжилгээ}

\subsection{Өрсөлдөгчийн шинжилгээ}

Одоогоор дотоодын зах зээлд NFC-д суурилсан ширээний захиалгын систем
\cite{iso14443, nfcforum} гэсэн бүрэн эрхт шийдэл алга. Өрсөлдөөн нь голдуу QR код дээр суурилсан маягийн цэс, өөрөө үйлчлэх терминал (kiosk), байгууллага өөрөө
зориулалтын аппликешнээр захиалга боловсруулсан хэлбэрээр явагдаж байна.
Харьцуулалтыг Хүснэгт~\ref{tab:ch2-ursulduuch}~—т харуулав.

\begin{table}[h]
  \centering
  \caption{Өрсөлдөгч шийдлүүдийн харьцуулалт}
  \label{tab:ch2-ursulduuch}
  \small
  \begin{tabular}{@{}p{3.0cm}p{3.0cm}p{3.0cm}p{3.0cm}p{2.5cm}@{}}
    \toprule
    \textbf{Шалгуур} & \textbf{NFC (энэ төсөл)} & \textbf{QR код цэс} &
    \textbf{Kiosk терминал} & \textbf{Өөрийн апп} \\ \midrule
    Эхлүүлэх алхам & Нэг алхам (стикер хүртэх) & Камер нээж скан хийх &
    Терминал руу очих & Апп татах, бүртгүүлэх \\
    Ширээ таних & Автоматаар, өндөр нарийвчлалтай & Кодын чанараас
    хамаарна & Гар аргаар сонгоно & Гар аргаар сонгоно \\
    Апп татах шаардлага & Байхгүй (PWA) & Байхгүй & Байхгүй & Шаардлагатай \\
    Эдэлж хэрэглэдэг & Шошго, элэгдэл бага & Код арчигдаж, бүдгэрнэ &
    Машин элэгдэнэ & --- \\
    Нэгжийн зардал & Бага (ширеэ бүр 1 стикер) & Бага & Өндөр (төхөөрөмж) &
    Хөгжүүлэлт, маркетинг өндөр \\
    Гадна орчинд & Тэсвэртэй & Тэсвэртэй & Хязгаартай & --- \\
    \bottomrule
  \end{tabular}
\end{table}

\subsection{Хэрэглэгч, зах зээлийн шинжилгээ}

\textbf{Хэрэглэгчийн бүлэг:}

\begin{itemize}
  \item \textbf{Үндсэн:} 18 -- 40 насны, гар утас болон цахим төлбөрийг
        өдөр тутамдаа ашигладаг, цагийн баримтлалтай хотын оршин суугчид
        (оюутан, цалинтай ажилтан, мэргэжилтнүүд).
  \item \textbf{Хоёрдогч:} Багийн ажил, уулзалт хийх хүмүүс, гадаадын
        жуулчид (англи хэл дэмжлэгтэйгээр).
  \item \textbf{B2B хэрэглэгч:} Салбар бүхий кофе шопууд, эмийн болон
        хүнсний салбарын сүлжээ дэлгүүрүүдийн хоол, ундааны
        үйлчилгээний хэсэг.
\end{itemize}

\textbf{Зах зээлийн үнэлгээ:} Улаанбаатар хотод 100 гаруй тооны кофе шоп
байрлах бөгөөд кофе шоп тус бүр 8 -- 30 ширээтэй. Эхний жилд 10 -- 20
кофе шопт (ширеэний тоогоор 200 -- 500 ширээ) бүтээгдэхүүнээ
борлуулах нь төслийн эхний зорилт бөгөөд салбарын нийт ширээний
тоогоор 10 -- 15\% хүрэх хувь хэмжээ юм.

\subsection{SWOT шинжилгээ}

\begin{table}[h]
  \centering
  \caption{SWOT шинжилгээ}
  \label{tab:ch2-swot}
  \small
  \begin{tabular}{@{}p{7.0cm}p{7.0cm}@{}}
    \toprule
    \textbf{Давуу тал (S)} & \textbf{Сул тал (W)} \\ \midrule
    NFC нь нэг алхамаар ширээг таньж, цэс нээдэг; апп таталгүйгээр
    ашиглах боломжтой; орчны нөлөөнд тэсвэртэй; бүрэн ажиллагааны
    зардлын хэмнэлт; дотоодын зах зээлд анхдагч давуу байдал &
    Бренд танигдаагүй байдал; NFC-ийн талаарх ойлголт хэрэглэгчид
    бага; хөгжүүлэлтийн хугацаанд хүний нөөцийн хязгаарлалт; эхний
    үеийн санхүүгийн хязгаарлалт \\ \midrule
    \textbf{Боломж (O)} & \textbf{Аюул (T)} \\ \midrule
    Цахим төлбөрийн хэрэглээ өсөж байна; кофе шопын салбар өсөж,
    цифржилтийн эрэлт нэмэгдэж байна; уламжлалт бус салбарын
    хамтран ажиллах боломж; гадаад зах зээлд экспортоор гарах
    боломж &
    QR кодын хямд, хурдан шийдлүүдийн өрсөлдөө; томоохон
    компаниудын ижил төстэй бүтээгдэхүүн гаргах эрсдэл; төлбөрийн
    үйлчилгээ үзүүлэгчдийн комиссийн ханшийн өөрчлөлт; хууль
    эрхзүйн шинэ зохицуулалт гарах \\ \bottomrule
  \end{tabular}
\end{table}

\noindent\textbf{Дүгнэлт:} SWOT-ын шинжилгээгээр төсөл нь дотоодын зах
зээлд анхдагч бүтээгдэхүүн болох давуу байдлыг ашиглан, QR кодоос
илүү хурдан, найдвартай туршлагаар өрсөлдөх боломжтой нь тодорхой
болов. Сул талыг нөхөхийн тулд маркетингийн чиглэлээр NFC-ийн ашиг
тусыг тодорхой, хялбар тайлбарласан сурталчилгаа, анхдагч хэрэглэгч
кофе шопуудтай гэрээ хийх төлөвлөгөөг тусгасан.
